% !TeX root = ../Tex/main.tex

\section{Models}

We revisit the four puzzels raised by the animal examples in \autoref{sub:problem} (\nameref{sub:problem}):

\begin{enumerate}
  \item Why does the animal initiates attack and then provides assistance? Isn't it contradictory?
  \item Why does the animal initiates attack only in one of the two instances?
  \item What is rational? To attack and regret that decision, or not attack at all? With respect to what criteria do we define which is the best strategy?
  \item Which is the best strategy in evolutionary terms?
\end{enumerate}

The following section draws on two classical thoeries of games with the aim of answering the four puzzels. The explanations improve upon these theories and support new idea.

\subsection{Model \(\text{n}^{{\small{0}}}1\): Surprise as Shock} \label{sub:model1}
The first two classical theories we can appeal to are \textit{deterrence} and \textit{bargaining}. \textcite{reiterExploringBargainingModel2003} summarises these positions arguing that: "Classic deterrence theory predicts that as an imbalance of power between two states grows, war becomes more probable. The bargaining model, \textit{on the other hand}, argues that it is not an imbalance of power that causes war, but rather disagreement over the balance of power (and over the likely outcome of war)"\textit{\small{(Emphasis added.)}}

\vspace*{1em}
In the light of formal representations of conflic, the first question can be rephrased as follows: can we depict \nameref{sub:case1} as a deterrence game and \nameref{sub:case2}  as a bargaining game?

\vspace*{1em}
Three doubts naturally arise as obstacles to the adoption of this view.
First, there is no intelligible reason why the animal acts through deterrence in one occasion and thorugh bargaining in the other. Secondly, we desperately need a justification for our distiction. Otherwise, We would incur in the unplesaint situiation of modelling two identical interactions through different formalisation without any specific reason. Supposing, we are able to find thoeritic support for these puzzles, we would still need to explain why the animal regrets the attack. Therefore, the third doubt arises because it appears unlikely to explain why the animal regrets the attack by making a distinction between deterrence and bargaining.

\vspace*{1em}
The decisive, and \textit{omitted}, variable could be the element of surprise.
\textcite{Fearon_1995} argues that, \textit{if we assume that conflicts are costly}, no war is possible when the outcome is clear to both sides. However, if we add surprise to \textcite{Fearon_1995}'s framework, we can derive an equilibrium.
The two versions of the elephant-human encounter are similar in strategic structure. In the tourist scenario, there is no surprise: the elephant is aware that humans might approach and disturb it, and it knows where to find them.

The situation involving the shepherd, however, contains the element of surprise, which may justify a different behavioural response.

\vspace*{1em}
In order to respect time and pages constraint, the main idea presented here cannot be further developed. The reason behind this choice is to prioritise the model that is introduced in the next section.

%However, before jumping to the model, we must ourselves what does empathy have to do with the first case. We may argue that empathy weakens elephant's signal. Yet empathy also functions as a force that regulates the intensity of the signal itself.
%Is it therefore reasonable to suggest that, in the second case, empathy may also be the mechanism that modulates the variability of the signal?
%What is the relationship between these two forces?
%Now that the cause has been established, we can turn to the effects. We observe two primary effects: that of surprise and that of empathy.

\subsection{Model \(\text{n}^{{\small{0}}}2\): Vote Buying}

\subsubsection{Back to Groseclose and Snyder (1996)}
This section develops I propose a reformulation of a well-known formalisation by \textcite{grosecloseBuyingSupermajorities1996}. I intend to draw on the equilibrium the authors derive, to support the argument developed in the DAG by providing an additional illustration. In doing so, I aim to improve upon the previous model (See \autoref{sub:model1}: \nameref{sub:model1}).

\vspace*{1em}
In their article, \textcite{grosecloseBuyingSupermajorities1996} explain why, in the US Congress, the number of votes often exceedes the minimum required to pass a bill. They argued that this is due to the strategic interaction between lobbying groups competing for support on different bills. These groups \textit{bribe} members of Congress to convince them. 

\vspace*{1em}
The central idea in the paper is illustrated by:

\vspace{1em}
\textit{"The legislature contains seven members, all indifferent between the two proposals, that is, $v(i) = 0$ for all $i$. Since B will move second and attack the weakest part of A's coalition, A will offer the same bribe, $a$, to all legislators he bribes. Let $m + 4 \ge 0$ be the size of the coalition A bribes. If $m = 0$, then B can defeat $x$ by paying $a + \varepsilon$ to one member of A's coalition. Thus, A must set $a \ge W_B$ to win, and A's total bribes are then $4W_B$. If $m = 1$, then B must pay $a + \varepsilon$ to two members of A's coalition to defeat $x$. Thus, A must set $a \ge W_B/2$ to win, and his total bribes are then $5W_B/2$. This is less than $4W_B$, so A is better off buying a supermajority of size 5 than a minimal winning coalition. If $m = 2$, then B must pay $a + \varepsilon$ to three members of A's coalition, so A must set $a \ge W_B/3$ to win, and his total bribes are then $2W_B$. In fact, A minimizes his total payments by bribing all seven legislators. Here, $m = 3$, and A's total bribes are $7W_B/4$. Clearly, there is an equilibrium where $x$ wins if and only if $W_A \ge 7W_B/4$. If $W_A < 7W_B/4$, then there are only equilibria in which no bribes are paid and $s$ wins."} \parencite{grosecloseBuyingSupermajorities1996}


\subsubsection{Representation through Vote Buying}
Consider a council composed of several members. These members are non-strategic players whose votes can be influenced. There are two active players, conceptualised as two groups, each seeking to "buy" the votes of the council members in order to determine the policy agenda. Each group possesses a budget, and bills are voted on by the council according to a pre-defined rule, assumed here to be the simple majority rule. This setup corresponds to the classical \textit{vote-buying game} as formalised by \parencite{grosecloseBuyingSupermajorities1996}.

\vspace*{1em}
The objective of this analysis is to adapt this framework to the so-called \textit{elephant-shaped} case.  
This adaptation allows the original model's equilibrium structure to be embedded in a cognitive-emotional set-up.


\begin{itemize}
    \item Let \( K \) denote the total number of \textit{thoughts}.
    \item The two players represent \textit{opposing emotions} that govern the internal dynamics of the elephant's cognition.  
    "Setting the agenda" refers to determining the behaviour.
    \item Each player's \textit{budget} corresponds to the \textit{intensity} of its emotion.  
    A higher emotional intensity enables the emotion to "purchase" (secure) a larger number of thoughts.
\end{itemize}

Notation:

\begin{itemize}
  \item[] Players: \(N=\{X,Y\}\) or \(N=\{Emotion 1, Emotion 2\}\).
  \item[] Histories: \(H= \big \{\varnothing,\{x\},\{x,y\}\big \}\). Wwere \(x\) and \(y\) indicate payment vectors in the baseline setting. Emotions use units of payment (emotional intensity) to corrupt a thought (or a judgement).
  \item[] Player function: \(p(\varnothing)=X\);\quad \(p(x)=Y\)
  \item[] Actions: \(A(\varnothing)=\{x\}\);\quad \(A(x)=\{y\}\) \(\forall\,x\).
  \item[] Each player \textit{receives} \(V_i\) if the elephant acts accordingly.
  \item[] Each player loses \(\sum_{j=1}^{k} x_{j}\) for each unit of intensity spent.
\end{itemize}

%\vspace{1em}
This adapted framework enables a formal representation of conflict as a strategic interaction between emotions competing for cognitive control. Each emotion allocates its limited intensity across thoughts, attempting to secure a majority sufficient to determine the elephant's conduct choice. The resulting equilibrium can be interpreted as a stable emotional configuration in which neither emotion can improve its outcome given the other's allocation strategy. Let \( \mu \) denote the minimal number of thoughts required to \textit{set the agenda}, that is, to achieve behavioural dominance of one emotion over the other. Then:
\[
\mu=\dfrac{k+1}{2}
\]



